\clearpage
\pdfbookmark[1]{Fourier. 1833 Februar 25.}{anzeige_fourier_128_130}
\SourceLeaf{128}
\begin{center}
\rule{0.55\linewidth}{0.4pt}\par Göttingische gelehrte Anzeigen.\quad 1833 Februar 25.\par \rule{0.55\linewidth}{0.4pt}
\end{center}
\textit{Analyse des équations déterminées par} \textsc{Fourier}.\quad\textit{Première partie.}\quad Paris 1831. \textit{Chez Firmin Didot frères.}\quad (XXIV und 258 S. in Quart.)


Nach dem Entwurf des Verfs. sollte dieses Werk aus zwei Abtheilungen bestehen, und die erstere ausser einer allgemeinen Uebersicht des Ganzen die beiden ersten Abschnitte, die zweite Abtheilung hingegen die fünf übrigen Abschnitte enthalten. Der Verf. wurde den Wissenschaften durch den Tod entrissen, als der Druck eben angefangen war. Indessen fand sich die ganze erste Abtheilung, die wir hier durch die Besorgung des Hrn. \textsc{Navier} erhalten, so gut wie vollständig ausgearbeitet vor, während der Rest von der Vollendung noch weit entfernt war: man muss dies um so mehr bedauern, da aus der allgemeinen Uebersicht hervorgeht, dass manche interessante Untersuchungen dafür bestimmt waren. Hier müssen wir, mit Uebergehung dessen, was wir zu erwarten gehabt hätten, wenn dem Verf. ein längeres Leben zu Theil geworden wäre, uns auf die Anzeige von dem einschränken, was wir wirklich erhalten haben.

Der Zweck des Werks ist, die numerische Auflösung der bestimmten algebraischen Gleichungen in einen ganz sichern und möglichst einfachen Gang zu bringen, namentlich die Anzahl der reellen und imaginären Wurzeln mit völliger Bestimmtheit auszumitteln, jene einzeln in feste Grenzen einzuschliessen und die stufenweise zu beliebiger Schärfe zu führende Annäherung zu ihren Werthen ganz methodisch anzuordnen. Die Hauptgrundlage des Ganzen bildet ein dem Verf.
\SourceLeaf{129}
eigenthümliches Theorem, welches als eine glückliche Generalisirung von \textsc{Descartes} Lehrsatz (gewöhnlich nach \textsc{Harriot} benannt) zu betrachten ist, und, dem Wesen nach, in Folgendem besteht. Es sei $X$ eine geordnete ganze algebraische Function von $x$ (mit lauter reellen Zahlencoëfficienten); ferner $a,a'$ zwei beliebige bestimmte ungleiche reelle Grössen, und zwar $a'-a$ positiv; durch die Substitutionen $x=y+a,\ x=y+a'$ gehe $X$ in die gleichfalls geordneten Functionen $Y,Y'$ über, in welchen die Coëfficienten resp. $g,g'$ Zeichenfolgen enthalten: der ausnahmliche Fall, wo einer oder mehrere Coëfficienten verschwinden, wird hier der Kürze wegen bei Seite gesetzt. Dies vorausgesetzt, kann die Gleichung $X=0$ innerhalb der Grenzen $x=a,\ x=a'$ nicht mehr als $g'-g$ reelle Wurzeln enthalten, oder noch bestimmter, wenn die Anzahl der reellen Wurzeln zwischen jenen Grenzen $=\lambda$ ist, so ist allemal $g'-g-\lambda$ entweder $=0$, oder eine gerade positive Zahl. Ist demnach die Differenz $g'-g$, welche niemals negativ sein kann, $=0$, so gibt es zwischen jenen Grenzen gar keine reelle Wurzel; wird $g'-g=1$, so enthalten die Grenzen eine reelle Wurzel und nicht mehr; ist endlich $g'-g=2$. so bleibt einstweilen noch ungewiss, ob zwei reelle Wurzeln zwischen den Grenzen liegen oder gar keine. Anstatt zweier Grenzen $a,a'$ kann man eine grössere Anzahl auf ganz ähnliche Art behandeln, und solche so wählen, dass erstlich der kleinsten eine Function $Y$ mit blossen Zeichenwechseln, der grössten eine mit blossen Zeichenfolgen entspricht, also wenn die entsprechenden Zahlen der Zeichenfolgen der Reihe nach mit $g,g',g'',g'''$ u.\,s.\,f. bezeichnet werden, die erste, dieser Zahlen $=0$, die letzte der Ordnungszahl der Gleichung gleich wird; und zweitens, dass sämmtliche einzelne Differenzen $g'-g,\ g''-g',\ g'''-g''$ u.\,s.\,w. nur entweder $=1$ oder $=2$ werden\quad Dadurch werden dann sämmtliche reelle Wurzeln dergestalt zwischen Grenzen eingeschlossen, dass in jedem einzelnen Intervall entweder eine liegen muss, oder zwei liegen können. Wie im letztern Fall auf ganz methodische Art, nöthigenfalls durch weitere Verengung der Grenzen zur Entscheidung gebracht wird, ob die zwei reellen Wurzeln wirklich vorhanden sind, oder fehlen, kann hier des beschränkten Raumes wegen nicht näher ausgeführt werden. Wir bemerken nur, dass so oft dieser letzte Fall eintritt, es allemal innerhalb solcher Grenzen einen Zwischenwerth gibt, für welchen in der Function $Y$ ein Coëfficient vor dem letzten ausfällt, während der vorhergehende und folgende gleiche Zeichen haben. \textsc{Fourier} nennt solche Stellen kritische: jede solche kritische
\SourceLeaf{130}
Stelle bedingt demnach das Fehlen von zwei reellen Wurzeln: wenn aber \textsc{Fourier} sich zugleich so ausdrückt, dass jedesmal zwei solche ausfallende reelle Wurzeln \textit{imaginär} werden, so können wir diesen Ausdruck nicht ganz billigen, da er leicht zu einer Misdeutung Veranlassung geben könnte. In der That ist es zwar wahr, dass die Gleichung $X=0$ \textit{zusammen gezählt} genau so viele Paare imaginärer Wurzeln enthält, als solche Ausfälle oder kritische Stellen vorkommen: allein die Werthe aller imaginären Wurzeln sind an sich eben so bestimmte Grössen wie die reellen, und jener Ausdruck kann daher leicht so gedeutet werden, als ob jeder \textit{bestimmten} Lücke ein \textit{bestimmtes} Paar imaginärer Wurzeln angehörte, was jedoch nicht nur von \textsc{Fourier} \textit{nicht} nachgewiesen ist, sondern so lange, als tiefer eindringende Untersuchungen diesen interessanten Punkt noch nicht in helles Licht gesetzt haben, zweifelhaft bleiben muss. Uebrigens soll hiermit nicht gesagt werden, dass \textsc{Fourier} selbst den Ausdruck so verstanden habe; wir möchten eher das Gegentheil annehmen, und fast vermuthen, dass er über das Dasein oder Nichtdasein eines solchen bestimmten Zusammenhanges ungewiss geblieben, und absichtlich einer offenen Erklärung über den verfänglichen Ausdruck ausgewichen sei. Ueberhaupt hat \textsc{Fourier} in das Wesen und die Berechnung der imaginären Wurzeln in diesem Werke sich gar nicht eingelassen, und es bleibt daher noch ein weites Feld zu bearbeiten übrig.\par
Einem so gewandten Grössenforscher, wie \textsc{Fourier} war, konnte es. einmal im Besitz jenes schönen Lehrsatzes --- und nach den von Herrn \textsc{Navier} mitgetheilten Notizen ist jener in diesem Besitz schon seit sehr langer Zeit gewesen --- nicht schwer fallen, auf demselben die Anordnung der Technik der numerischen Auflösung der Gleichungen zu begründen, und diese Entwickelung ist mit grosser Vollständigkeit und Ausführlichkeit gegeben. Geübtere Leser möchten vielleicht eine etwas gedrängtere Darstellung und die Wegschneidung mancher Wiederholungen vorziehen, dem weniger geübten werden die vielen gut gewählten und ausführlich behandelten Beispiele willkommen sein. Jedenfalls sichert dieses Werk \textsc{Fourier}'s Namen auch in diesem Theile der Grössenlehre einen ehrenvollen Platz, den er in andern schon längst behauptet.\par
\par\noindent\hfill\rule{0.30\linewidth}{0.4pt}\hfill\mbox{}\par

